\FloatBarrier
\section{Modelling of cement hydration and its reliance on \acs*{SEM} images}

% aus dem ct/fib-sem-antrag
As the demands on building materials increase, so do the requirements for characterising and modelling the underlying material systems.
There are advanced approaches to modelling and simulating the strength, shape retention and durability characteristics of concretes \cite{bishnoi.2009, kumar.2017, zhou.2019, ouzia.2019}.
Due to the high complexity, multi-scale heterogeneity and temporal changes, this is often only possible in a highly simplified model and on a very limited scale.

The microstructure of concretes was simulated in a variety of models at different scales \cite{bishnoi.2009, kumar.2017, zhou.2019, ouzia.2019}.
The mesostructure (between atomistic and micro scale) of e.g. \ac{CSH} phases and the porosity was also modelled \cite{ioannidou.2016, yu.2016}.
Atomistic modelling is used to fundamentally simulate the structural composition \cite{kumar.2017,mishra.2017,kunhimohamed.2020,casar.2023}.
Finally, there are modern micromechanical models that attempt to close the transition from meso- to macro-modelling \cite{bernard.2003, hlobil.2016, gobel.2018}.

The hydration of cement over time can also be simulated using the sheet growth model in 2D \cite{etzold.2014} and 3D \cite{nguyen-tuan.2019,nguyen-tuan.2023} (see \zcref{fig:sheet_growth}).

% from a talk of long from the ICCC 2019, slide 14/18 but not in the paper
\begin{figure*}[htb]
    \centering

    \begin{subfigure}[t]{0.32\linewidth}
        \includegraphics[trim={186 0 267 44},clip,width=\textwidth]{literature/sheet_grow_anim/sg01.jpg}
    \end{subfigure}%
    \hfill%
    \ifdefined\PRINTABLE
        \begin{subfigure}[t]{0.32\linewidth}
            \includegraphics[trim={186 0 267 44},clip,width=\textwidth]{literature/sheet_grow_anim/sg02.jpg}
        \end{subfigure}
    \else
        \animategraphics[trim={186 0 267 44},clip, palindrome,autoplay,poster=02,width=.32\linewidth]{3}{literature/sheet_grow_anim/sg}{01}{04}
    \fi
    \hfill%
    \begin{subfigure}[t]{0.32\linewidth}
        \includegraphics[trim={186 0 267 44},clip,width=\textwidth]{literature/sheet_grow_anim/sg04.jpg}
    \end{subfigure}

    \caption{
        Visualisation of the growth of \ac{CSH} in a representative domain of \qtyproduct{.5 x .5 x .5}{\micro\meter\cubed} between \qty{15}{\minute} and \qty{4}{\hour}, based on the sheet growth model \cite{nguyen-tuan.2019,nguyen-tuan.2023}. % in the second paper 300x300x2200nm is used due to some restrictions.
    }
    \label{fig:sheet_growth}
\end{figure*}
For larger domains, \textcite{nguyen-tuan.2024} utilised the level set-based method, allowing simulation of the growth of low- and high-density \ac{CSH} in 2D based on a segmented image of polished alite after \qty{1}{\day} of hydration (see \zcref{fig:level_set}).

\begin{figure*}[htb]
    \centering

    \begin{subfigure}[t]{0.32\linewidth}
        \includegraphics[width=\textwidth]{literature/level_set_anim/ls01.jpg}
    \end{subfigure}%
    \hfill%
    \begin{subfigure}[t]{0.32\linewidth}
        \ifdefined\PRINTABLE
            \includegraphics[width=\textwidth]{literature/level_set_anim/ls04.jpg}
        \else
            \animategraphics[palindrome,autoplay,poster=04,width=\textwidth]{6}{literature/level_set_anim/ls}{01}{09}
        \fi
    \end{subfigure}%
    \hfill%
    \begin{subfigure}[t]{0.32\linewidth}
        \includegraphics[width=\textwidth]{literature/level_set_anim/ls07.jpg}
    \end{subfigure}

    \caption{Visualisation of the growth of \ac{CSH} (purple) in the low and high density variant on top of \ac{C3S} (yellow) by \textcite{nguyen-tuan.2024}, based on the level set method. The red line indicates the surface of the unhydrated \ac{C3S} particle.}
    \label{fig:level_set}
\end{figure*}
%{\color{red} % erstmal nur so dahin gequatscht
\textsc{CEMHYD3D} \cite{bentz.2005} is a programme to simulate the hydration of cement in three dimensions using a voxel-based model.
Models like these (\cite{bentz.2005, nguyen-tuan.2024}) allow simulation of the development of the phase composition over time, which can be validated using various experiments.
3D-based simulations inherently have a better chance to closely represent reality, since 2D simulations are always missing a lot of information, which can cause distorted results (see \zcref{subsec:sectioning-effect}).
However, 2D simulations are much faster and allow simulation of larger areas, reflecting a more representative area, which is very important due to the heterogeneity of cement (see \zcref{sec:2D-representativeness}).
% simulationspaper der hydratation \cite{hlobil.2022} which was compared to volumetric \ac{mCT} data of differently aged cement.
% CEMHYD3D -> Fig 3 anteile der phasen vgl mit Messungen
% Literatur aus Long übernehmen und Modelle erwähnen
%}

Nevertheless, the basis for all realistic modelling is a qualitative and quantitative experimental characterisation of the corresponding structures \cite{nguyen-tuan.2019, gobel.2018a}.


% Handlungsanweise - wie kann man REM mit Methoden kombinieren um Baustoffe zu beschreiben - welche Techniken gibt es?
% Einfache Tools, weil sonst machts keener

\pdfcomment[icon=Note, color=orange]{% TODO: note
    Noch einarbeiten:
    HYMOSTRUC, HYMOSTRUC3D - vanbreugel.1995, ye.2007, liu.2014, gao.2014, qian.2017, gao.2019
}

